\chapter{Преглед на предметната област}

\section{Основни дефиниции, концепции и алгоритми}
Електронната таблица се моделира като правоъгълна матрица в оперативната памет. Основните градивни единици са клетките, които могат да бъдат празни или да съдържат: числа, текст (в кавички) или формули (започващи с \texttt{=}). За пресмятане на формулите и извличане на стойности от други клетки се използва рекурсивен алгоритъм.

\section{Проблеми и сложност}
Основно предизвикателство е как формулите с препратки да извличат нужните стойности, без да се създава силна зависимост към цялостната структура на таблицата.

Друг съществен проблем възниква при зареждането на данните. Динамичното изчисляване на формули не може да се извършва линейно по време на четенето от файл, тъй като те могат да реферират към клетки, които все още не са заредени в паметта.

\section{Методи за решаване на проблемите}
За справяне с дефинираните алгоритмични предизвикателства са приложени следните подходи:

\begin{itemize}
  \item \textbf{Взаимодействие и достъп до данни:} За да се избегне подаването на директен достъп до целия масив с данни, се използва механизъм за индиректно извличане на стойности. По този начин формулите получават нужните данни в реално време, без да се нарушава капсулацията на таблицата.

  \item \textbf{Отложено изчисляване на формули:} Първо цялата таблица се зарежда напълно в паметта, а след това за извличане на липсващи стойности от реферирани клетки се използва рекурсия. За да се предотврати безкраен цикъл при взаимнозависими клетки (циклични препратки), преди всяко извикване временно се задава статус на активно изчисляване. При повторно засичане на такъв статус, алгоритъмът разпознава зациклянето и безопасно прекъсва операцията.
\end{itemize}

\section{Потребителски и качествени изисквания}

\textbf{Функционални изисквания (Потребителски):}
Приложението предоставя конзолен интерфейс (CLI). Таблиците се зареждат от текстов файл (стойности, разделени със запетаи). Потребителят разполага със следните команди:
\begin{itemize}
  \item Файлови: \texttt{open <file>}, \texttt{close}, \texttt{save}, \texttt{saveas <file>}, \texttt{help}, \texttt{exit}.
  \item Операционни: \texttt{edit <address> <value>} (промяна на клетка) и \texttt{print} (отпечатване с автоматично подравняване).
\end{itemize}

\textbf{Качествени изисквания (Нефункционални):}
Системата трябва да осигурява висока стабилност при невалидни данни – да конвертира празни и текстови клетки до 0 при изчисление, и да индикира грешки (напр. деление на нула) без срив на програмата. Приложението трябва да е устойчиво на неточно въвеждане, като автоматично игнорира излишни интервали около разделителите.